% Chapter 6: Error Correcting Codes
% Part II: Information

\chapter{Error Correcting Codes}
\label{ch:error_correction}

%======================================================================
% SECTION 1: THE QUESTION
%======================================================================
\section{The Question: What Problem Was Humanity Trying to Solve?}
\label{sec:ec:question}

Shannon's 1948 Channel Coding Theorem proved that \emph{reliable} communication at rates up to channel capacity $C$ is possible — but it was an \emph{existence proof} using random codes. It did not tell engineers \emph{how} to build practical encoders/decoders.

How can we add redundancy to a message so that errors introduced by a noisy channel can be detected and corrected — efficiently, with practical encoding and decoding algorithms?

The problem: Noise is inevitable. Thermal noise, cosmic rays, crosstalk, fading, interference — every physical channel corrupts bits. Naive repetition (send each bit 3 times, take majority) works but is terribly inefficient (rate 1/3). We need codes that approach capacity with feasible complexity.

Every communication system faces the same trilemma: power, bandwidth, and reliability. Increasing transmitter power helps but is expensive and constrained by physics and regulation. Increasing bandwidth trades one scarce resource for another. The only free parameter is coding: can we add redundancy in a clever way that provides reliability without sacrificing throughput?

Engineers before Shannon believed that to communicate reliably over a noisy channel, your only options were to shout louder (increase power) or repeat yourself more often (decrease rate). The idea that you could add redundancy in a structured way and \emph{simultaneously} maintain high rate and high reliability seemed like a contradiction in terms.

Can we achieve Shannon's promised rates with \emph{structured} codes that have \emph{efficient} (polynomial-time) encoding and decoding algorithms? Or is the gap between theory and practice fundamental?

The stakes could not be higher. Every system that stores or transmits digital data — from the memory in your phone to the satellite transmitting images from Mars — depends on error correction. The difference between a functioning system and a broken one is often a few dB of coding gain. The search for practical codes that approach the Shannon limit is one of the great engineering quests of the information age.

Shannon proved that for any channel with capacity $C$, there exist codes at any rate $R < C$ with arbitrarily low error probability. But the proof was non-constructive — it used random codes and averaging arguments. This chapter is the story of the 60-year quest to make Shannon's promise concrete.

%======================================================================
% SECTION 2: WHY IT SEEMED IMPOSSIBLE
%======================================================================
\section{Why It Seemed Impossible: What Barriers Blocked Progress}
\label{sec:ec:impossible}

Before Shannon, the prevailing intuition was:

\begin{enumerate}
\item \textbf{Redundancy $\to$ Lower Rate:} Adding redundancy \emph{reduces} information rate. To correct errors, you must send extra bits. The more errors you want to correct, the more redundancy you need. Rate $\to$ 0 as reliability $\to$ 1.

\item \textbf{The ``Noise Barrier'':} Engineers believed there was a fundamental tradeoff: you can have \emph{either} high rate \emph{or} high reliability, but not both. The ``error floor'' was considered unavoidable.

\item \textbf{Decoding Complexity:} Even if good codes existed, \emph{finding} the closest codeword to a received word (Maximum Likelihood Decoding) seemed to require exponential search over all codewords. For a $(n,k)$ code, $2^k$ codewords — impossible for large $k$.
\end{enumerate}

Each of these barriers rested on hidden assumptions. The first assumed that all redundancy is repetition, ignoring the possibility of \emph{structured} algebraic redundancy. The second assumed worst-case error patterns, whereas Shannon showed we need only handle \emph{typical} errors. The third assumed brute-force search was the only decoding strategy, overlooking the power of algebraic structure and iterative algorithms.

For a $(1000, 500)$ code, there are $2^{500} \approx 10^{150}$ possible codewords — vastly more than the number of atoms in the observable universe. Brute-force maximum likelihood decoding is impossible at any scale. Yet modern LDPC decoders process codewords of this size in microseconds, using iterative message-passing algorithms that converge to near-ML performance.

``To communicate reliably, you must shout louder (increase power) or repeat yourself (decrease rate). There is no clever coding trick that beats physics.''

Shannon proved this intuition wrong — but his proof used \emph{random codes} with no structure, requiring exponential-time decoding. The challenge: find \emph{structured} codes with \emph{efficient} decoding.

The challenge of error-correcting codes mirrors the fundamental tension in communication theory: power, bandwidth, and complexity form a trilemma. Coding is the lever that lets us trade one for another. Before Shannon, engineers didn't know the lever existed.

%======================================================================
% SECTION 3: HISTORICAL CONTEXT
%======================================================================
\section{Historical Context: What Was Known at the Time}
\label{sec:ec:context}

\begin{itemize}
\item \textbf{1948:} Shannon proves capacity exists; random codes achieve it \cite{Shannon1948}.
\item \textbf{1950:} Hamming invents (7,4) Hamming code --- first perfect single-error-correcting code. Algebraic coding theory born \cite{Hamming1950}.
\item \textbf{1954--57:} Golay codes (perfect binary and ternary) \cite{Golay1949}.
\item \textbf{1959--60:} Bose-Chaudhuri-Hocquenghem (BCH) codes; Reed-Solomon codes (non-binary, burst-error correction) \cite{ReedSolomon1960}.
\item \textbf{1960:} Peterson: Decoding algorithm for BCH codes.
\item \textbf{1966:} Forney: Concatenated codes (inner + outer) $\to$ exponentially decreasing error probability with polynomial complexity \cite{Forney1966}.
\item \textbf{1967:} Gallager: LDPC codes (forgotten until 1990s) \cite{Gallager1963}.
\item \textbf{1973:} Viterbi algorithm (for convolutional codes) \cite{Viterbi1967}.
\item \textbf{1993:} Berrou, Glavieux, Thitimajshima: Turbo codes --- near capacity with iterative decoding \cite{Berrou1993}.
\item \textbf{1996:} MacKay, Neal; Richardson, Urbanke: Rediscovery of LDPC codes \cite{gallager1968}.
\item \textbf{2009:} Ar{\i}kan: Polar codes --- first \emph{provably} capacity-achieving for symmetric channels with explicit construction and $O(n \log n)$ decoding \cite{Arikan2009}.
\end{itemize}

Key figures:
\begin{itemize}
\item \textbf{Richard Hamming (1915--1998):} Bell Labs. Frustrated by computer crashes on weekends --- wanted automatic error correction \cite{Hamming1950}.
\item \textbf{Irving Reed \& Gustave Solomon (1960):} Polynomial-based codes over finite fields. Foundation of CDs, DVDs, QR codes, DSL, RAID \cite{ReedSolomon1960}.
\item \textbf{Robert Gallager (1931--2021):} MIT. PhD thesis (1963) invented LDPC codes and iterative message passing. Ignored for 30 years \cite{Gallager1963,gallager1968}.
\item \textbf{Claude Berrou (1951--):} Telecom Bretagne. Turbo codes (1993) shocked the community --- ``too good to be true'' \cite{Berrou1993}.
\item \textbf{Erdal Ar{\i}kan (1958--):} Bilkent. Polar codes (2009) --- channel polarization, deterministic construction \cite{Arikan2009}.
\end{itemize}

Richard Hamming joined Bell Labs in 1946, working on the Bell Model V relay computer — one of the largest electromechanical computers ever built. The machine ran complex calculations unattended over weekends, but parity-check circuits would detect errors and halt the machine. When Hamming arrived Monday morning, he would find the machine stopped and have to reload the program from scratch. Frustrated, he asked: ``If the machine can detect an error, why can't it locate and correct it?'' This question led directly to the first error-correcting codes.

Hamming's 1950 paper, \emph{Error Detecting and Error Correcting Codes}, founded the field of algebraic coding theory. His code was the first to show that redundancy could be added in a structured way such that errors were not just detected but automatically corrected. Every DRAM module in the world still uses a variant of Hamming's original code for single-error correction. Hamming later contributed to numerical analysis, operating systems, and the famous ``Hamming distance'' that bears his name. He once said: ``The purpose of computing is insight, not numbers.''

Robert Gallager's 1963 MIT PhD thesis, \emph{Low-Density Parity-Check Codes}, was decades ahead of its time. He proposed using sparse parity-check matrices — matrices with very few 1s per row and column — and decoding them with an iterative message-passing algorithm. The approach was elegant and theoretically well-motivated, but the computational power of the 1960s made practical implementation impossible. The thesis was published as a book by MIT Press but largely ignored.

Thirty years later, the rediscovery of LDPC codes by MacKay and Neal (1996) and independently by Richardson and Urbanke (2001) sparked a revolution. Gallager's iterative decoding principle turned out to be not just a clever trick but a fundamental algorithmic paradigm — belief propagation on graphical models. Today, LDPC codes are the standard for WiFi (802.11n/ac/ax), 5G data channels, satellite communications, and SSD storage. Gallager also made foundational contributions to network information theory, including the capacity region of multiple-access channels and the celebrated ``Gallager bound.'' He received the Shannon Award in 1983 and passed away in 2021, having seen his forgotten thesis transform global communications.

At the 1993 IEEE International Conference on Communications in Geneva, Claude Berrou, Alain Glavieux, and Punya Thitimajshima presented a paper that shook the coding world: \emph{Near Shannon Limit Error-Correcting Coding and Decoding: Turbo-Codes}. Their codes performed within 0.5 dB of the Shannon limit — a performance previously thought impossible for any practical code. The audience was skeptical; many believed the results were a mistake or a simulation error.

Berrou named them ``turbo'' codes after the turbocharger in combustion engines, because of the feedback loop between two component decoders. The key insight was \emph{iterative decoding with extrinsic information exchange}: each decoder produces soft estimates of each bit, and the other decoder uses these estimates as prior information. The loop iterates until convergence. This principle — local message passing with feedback — turned out to be deeply connected to Gallager's earlier work and to belief propagation in probabilistic graphical models. Turbo codes were adopted in 3G and 4G cellular standards and remain one of the most impactful coding inventions of the 20th century.

Erdal Arıkan, a professor at Bilkent University in Ankara, Turkey, spent over a decade searching for a provably capacity-achieving code with explicit construction and efficient decoding. The problem had been open since Shannon's 1948 paper: could one write down a deterministic, constructive family of codes that provably achieve capacity with polynomial-time encoding and decoding?

His breakthrough came in 2008 with the concept of \emph{channel polarization}. Arıkan showed that by recursively combining and splitting independent copies of a binary-input memoryless channel, the channels ``polarize'': some become noiseless (capacity approaching 1) while others become completely noisy (capacity approaching 0). By transmitting information only through the noiseless channels, one achieves the symmetric capacity of the original channel. The construction was explicit, the encoding and decoding complexity was $O(N \log N)$, and the proof was rigorous. The coding community was stunned — a problem open for 60 years had been solved.

In 2016, 3GPP adopted polar codes for the 5G NR control channel, making them the first provably capacity-achieving code to be standardized. Arıkan received the 2018 Claude E. Shannon Award and the 2019 IEEE Richard W. Hamming Medal.

Shannon's Channel Coding Theorem (Chapter 5) is the \emph{existence proof}. This chapter is the \emph{constructive quest} to realize that promise.

%======================================================================
% SECTION 4: THE MENTAL LEAP
%======================================================================
\section{The Mental Leap: The Creative Insight That Changed Everything}
\label{sec:ec:leap}

There isn't just one leap — there are \emph{four} major leaps, each breaking a different barrier:

\emph{Algebraic Structure.} Represent codewords as vectors in $\mathbb{F}_2^n$ (or $\mathbb{F}_q^n$). Linear codes = subspaces. Encoding = matrix multiplication. Syndrome decoding = linear algebra. Errors become solvable equations.

\emph{Concatenation.} Don't try to build one perfect code. Build an \emph{inner} code for the channel (handles noise) and an \emph{outer} code for residual errors (handles bursts). The combination achieves exponentially low error with polynomial complexity.

\emph{Iterative Decoding on Sparse Graphs.} Long, random-like codes (LDPC, Turbo) decoded by \emph{message passing} on a Tanner graph. Local computations $\to$ global inference. ``Belief propagation'' approximates ML decoding.

\emph{Channel Polarization.} Take $N$ independent copies of a channel. Combine them recursively. Some become \emph{perfect} (capacity 1), others \emph{useless} (capacity 0). Send data only through perfect ones. Deterministic, explicit, $O(N \log N)$.

Each leap changed the \emph{design philosophy}: Algebra $\to$ Concatenation $\to$ Probabilistic/Graphical $\to$ Polarization/Deterministic.

Each leap built on the limitations of the previous. Hamming's algebra worked beautifully for short codes but could not scale to long blocklengths. Forney's concatenation achieved exponentially decaying error probability but could not reach the Shannon limit. Gallager and Berrou's iterative decoding reached capacity empirically but lacked a rigorous proof. Arıkan's polarization provided both the proof and the explicit construction — closing the 60-year gap between Shannon's existence theorem and a practical, provably capacity-achieving code.

%======================================================================
% SECTION 5: THE MATHEMATICS
%======================================================================
\section{The Mathematics: Rigorous Development of the Theory}
\label{sec:ec:math}

\subsection{Definitions}
\label{sec:ec:definitions}

\begin{definition}[Block Code]
An $(n, M, d)$ \emph{block code} over alphabet $\mathcal{X}$ is a set $\mathcal{C} \subseteq \mathcal{X}^n$ of $M$ codewords. Minimum Hamming distance $d = \min_{c \neq c'} d_H(c, c')$.
\end{definition}

\begin{definition}[Linear Code]
A \emph{linear} $[n, k, d]$ code over $\mathbb{F}_q$ is a $k$-dimensional subspace of $\mathbb{F}_q^n$. Generator matrix $G \in \mathbb{F}_q^{k \times n}$: $\mathcal{C} = \{ uG \mid u \in \mathbb{F}_q^k \}$. Parity-check matrix $H \in \mathbb{F}_q^{(n-k) \times n}$: $\mathcal{C} = \{ c \mid H c^T = 0 \}$.
\end{definition}

\begin{definition}[Syndrome Decoding]
Received $r = c + e$. Syndrome $s = H r^T = H e^T$. Depends only on error $e$. If $e$ is correctable, $s$ uniquely identifies $e$.
\end{definition}

\begin{definition}[Code Rate]
The \emph{rate} of an $[n,k,d]$ linear code is $R = k/n$. The rate measures how many information bits are transmitted per channel use. The fundamental tradeoff: higher rate means less redundancy and weaker error correction; lower rate means more redundancy and stronger correction. Shannon's theorem provides the fundamental limit: reliable communication is possible iff $R < C$, where $C$ is the channel capacity.
\end{definition}

\subsection{Hamming Codes}
\label{sec:ec:hamming}

\begin{theorem}[Hamming Codes]
For any $r \ge 2$, there exists a perfect $[2^r-1, 2^r-1-r, 3]$ binary linear code (single-error-correcting).
\end{theorem}

\begin{proof}[Construction]
$H$ has all non-zero binary $r$-tuples as columns. $n = 2^r-1$, $n-k = r$. Any two columns are linearly independent $\implies d \ge 3$. Sphere-packing bound met with equality $\implies$ perfect.
\end{proof}

\begin{example}[(7,4,3) Hamming Code]
$r=3$: $H$ has columns $001, 010, 011, 100, 101, 110, 111$. $G = [I_4 \mid P]$. Syndrome $s = H r^T$ gives binary representation of error position.
\end{example}

\begin{example}[(7,4) Hamming Code: Detailed Worked Example]
Take the parity-check matrix with columns in binary order:
\[
H = \begin{pmatrix}
0 & 0 & 0 & 1 & 1 & 1 & 1 \\
0 & 1 & 1 & 0 & 0 & 1 & 1 \\
1 & 0 & 1 & 0 & 1 & 0 & 1
\end{pmatrix}
\]
The generator matrix in systematic form is:
\[
G = \begin{pmatrix}
1 & 0 & 0 & 0 & 1 & 1 & 0 \\
0 & 1 & 0 & 0 & 1 & 0 & 1 \\
0 & 0 & 1 & 0 & 0 & 1 & 1 \\
0 & 0 & 0 & 1 & 1 & 1 & 1
\end{pmatrix}
\]
where $G = [I_4 \mid P]$ and $H = [P^T \mid I_3]$.

Encode message $m = (1,0,1,1)$:
\[
c = mG = 1\cdot r_1 + 0\cdot r_2 + 1\cdot r_3 + 1\cdot r_4
\]
\[
= (1,0,0,0,1,1,0) + (0,0,1,0,0,1,1) + (0,0,0,1,1,1,1)
\]
\[
= (1,0,1,1,0,1,0)
\]

The 16 codewords of the $(7,4)$ Hamming code are all linear combinations of the rows of $G$:
\[
\begin{array}{ll}
0000000,\ 0001011,\ 0010110,\ 0011101,\\
0100111,\ 0101100,\ 0110001,\ 0111010,\\
1000101,\ 1001110,\ 1010011,\ 1011000,\\
1100010,\ 1101001,\ 1110100,\ 1111111
\end{array}
\]

Now suppose the transmitted codeword $c = (1,0,1,1,0,1,0)$ is corrupted by a single bit error in position 3, giving received word $r = (1,0,0,1,0,1,0)$. Compute the syndrome:
\[
s = H r^T = \begin{pmatrix}0&0&0&1&1&1&1\\0&1&1&0&0&1&1\\1&0&1&0&1&0&1\end{pmatrix} \begin{pmatrix}1\\0\\0\\1\\0\\1\\0\end{pmatrix}
= \begin{pmatrix}0\\1\\1\end{pmatrix}
\]

The syndrome $s = (011)_2 = 3$ matches column 3 of $H$, indicating the error is in bit position 3. Flipping bit 3 recovers the original codeword $(1,0,1,1,0,1,0)$. This is the magic of the Hamming code: a single matrix multiplication locates the error exactly.

The standard array for decoding uses the syndrome to look up the error pattern:
\[
\begin{array}{c|c}
\text{Syndrome} & \text{Error pattern} \\ \hline
000 & 0000000 \\
001 & 0000001 \\
010 & 0000010 \\
011 & 0000100 \\
100 & 0001000 \\
101 & 0010000 \\
110 & 0100000 \\
111 & 1000000
\end{array}
\]
Each non-zero syndrome uniquely identifies a correctable error pattern (a weight-1 vector), confirming the code is perfect: every possible syndrome corresponds to either no error or a single-bit error.
\end{example}

\subsection{Reed-Solomon Codes}
\label{sec:ec:rs}

\begin{theorem}[Reed-Solomon Codes]
For any $n = q-1$, $k \le n$, there exists an $[n, k, n-k+1]$ code over $\mathbb{F}_q$ (maximum distance separable, MDS).
\end{theorem}

\begin{proof}[Construction]
Message $u = (u_0, \dots, u_{k-1}) \to$ polynomial $p(x) = \sum u_i x^i$. Codeword: $(p(\alpha), p(\alpha^2), \dots, p(\alpha^n))$ where $\alpha$ is primitive in $\mathbb{F}_q$. Degree $< k$ polynomial has $\le k-1$ roots $\implies$ at most $k-1$ zeros $\implies d \ge n-k+1$. Singleton bound $\implies$ MDS.
\end{proof}

\begin{example}[Reed-Solomon Encoding: Detailed Worked Example]
Consider a $(7,3)$ RS code over $\mathbb{F}_8$ with primitive polynomial $p(x) = x^3 + x + 1$. Let $\alpha$ be a primitive element. The field $\mathbb{F}_8$ has elements $\{0, 1, \alpha, \alpha^2, \alpha^3, \alpha^4, \alpha^5, \alpha^6\}$ where $\alpha^3 = \alpha + 1$, $\alpha^4 = \alpha^2 + \alpha$, etc.

Message: $m = (\alpha, 1, \alpha^2)$. Construct polynomial:
\[
p(x) = \alpha + 1\cdot x + \alpha^2 \cdot x^2
\]

Evaluate at $\alpha^0, \alpha^1, \dots, \alpha^6$:
\[
\begin{aligned}
c_0 &= p(1) = \alpha + 1 + \alpha^2 = \alpha^2 + \alpha + 1 = \alpha^5 \\
c_1 &= p(\alpha) = \alpha + \alpha + \alpha^4 = \alpha^4 \\
c_2 &= p(\alpha^2) = \alpha + \alpha^2 + \alpha^6 = \alpha^6 + \alpha^2 + \alpha \\
c_3 &= p(\alpha^3) = \alpha + \alpha^3 + \alpha^8 = \alpha + (\alpha+1) + \alpha \\
c_4 &= p(\alpha^4) = \alpha + \alpha^4 + \alpha^{10} \\
c_5 &= p(\alpha^5) = \alpha + \alpha^5 + \alpha^{12} \\
c_6 &= p(\alpha^6) = \alpha + \alpha^6 + \alpha^{14}
\end{aligned}
\]

Codeword: $(c_0, c_1, c_2, c_3, c_4, c_5, c_6) \in \mathbb{F}_8^7$.

Now introduce two symbol errors at positions 1 and 4. The received word is $r = (c_0, e_1, c_2, c_3, e_4, c_5, c_6)$ where $e_1 \neq c_1$ and $e_4 \neq c_4$.

\textbf{Step 1: Compute syndromes.} For a code with minimum distance $d = 5$ (correcting $t=2$ errors), compute $S_j = r(\alpha^j)$ for $j = 1, 2, 3, 4$:
\[
S_1 = r(\alpha),\quad S_2 = r(\alpha^2),\quad S_3 = r(\alpha^3),\quad S_4 = r(\alpha^4)
\]

Since $r(x) = p(x) + e(x)$ where $e(x)$ is the error polynomial with non-zero terms at the error positions, the syndromes depend only on the errors:
\[
S_j = e(\alpha^j) = \sum_{l=1}^2 e_{i_l} \alpha^{j \cdot i_l}
\]
where $i_1, i_2$ are the error positions and $e_{i_1}, e_{i_2}$ are the error values.

\textbf{Step 2: Find error locator polynomial.} The error locator polynomial $\Lambda(x) = (1 - \alpha^{i_1} x)(1 - \alpha^{i_2} x) = 1 + \Lambda_1 x + \Lambda_2 x^2$ has roots at $x = \alpha^{-i_1}, \alpha^{-i_2}$. The Berlekamp-Massey algorithm (Section~\ref{sec:ec:bch}) finds $\Lambda(x)$ from the syndromes. The coefficients satisfy:
\[
S_1 \Lambda_2 + S_2 \Lambda_1 + S_3 = 0,\quad
S_2 \Lambda_2 + S_3 \Lambda_1 + S_4 = 0
\]

\textbf{Step 3: Chien search.} Evaluate $\Lambda(\alpha^{-i})$ for $i = 0, 1, \dots, 6$. When $\Lambda(\alpha^{-i}) = 0$, position $i$ is in error. This gives $\{i_1, i_2\}$.

\textbf{Step 4: Forney's formula.} Compute error values:
\[
e_{i_l} = \frac{\omega(\alpha^{-i_l})}{\Lambda'(\alpha^{-i_l})}
\]
where $\omega(x) = \Lambda(x)(1 + S(x)) \bmod x^{d-1}$ is the error evaluator polynomial, and $\Lambda'$ is the formal derivative. Subtracting the error values at the located positions recovers the original codeword.

This four-step procedure (syndromes $\to$ BM $\to$ Chien $\to$ Forney) is the standard algebraic decoding pipeline for RS and BCH codes, and is implemented in billions of devices worldwide.
\end{example}

Burst errors (scratches on CDs, fading in wireless) corrupt consecutive bits. RS codes over $\mathbb{F}_{2^m}$ treat each symbol as $m$ bits $\to$ burst of $b$ bits affects $\le \lceil b/m \rceil$ symbols. RS corrects up to $\lfloor (n-k)/2 \rfloor$ symbol errors. This is why CDs, DVDs, QR codes, DSL, RAID-6 all use RS.

\subsection{BCH Codes}
\label{sec:ec:bch}

Binary codes designed via roots in extension fields. Narrow-sense BCH: designed distance $\delta$, roots $\alpha, \alpha^2, \dots, \alpha^{\delta-1}$. $d \ge \delta$. Efficient algebraic decoding (Peterson-Gorenstein-Zierler, Berlekamp-Massey).

\subsubsection*{Peterson-Gorenstein-Zierler Algorithm (1960)}
The original algebraic decoding algorithm for BCH codes, limited to small error counts. It works by solving a system of linear equations:
\begin{enumerate}
\item Compute syndromes $S_1, \dots, S_{2t}$.
\item Form the syndrome matrix $M = [S_{i+j-2}]$ for $i,j = 1,\dots,t$.
\item If $\det(M) \neq 0$, solve $M \Lambda = -S$ for the error locator polynomial coefficients.
\item Find roots of $\Lambda(x)$ via Chien search $\to$ error positions.
\item For non-binary BCH/RS: compute error values via Forney's formula.
\end{enumerate}
Complexity: $O(t^3)$ for the matrix inversion, $O(n)$ for Chien search.

\subsubsection*{Berlekamp-Massey Algorithm (1969)}
Finds the minimal linear feedback shift register (LFSR) generating a sequence. For BCH/RS decoding:
\begin{enumerate}
\item Compute syndromes $S_1, \dots, S_{2t}$.
\item Run BM to find error locator polynomial $\Lambda(x)$.
\item Find roots (Chien search) $\to$ error positions.
\item Forney's formula $\to$ error values.
\end{enumerate}
Complexity: $O(n^2)$ or $O(n \log^2 n)$ with FFT.

The Berlekamp-Massey algorithm is more efficient than PGZ for large $t$ because it avoids explicit matrix inversion. It iteratively constructs the minimal polynomial that generates the syndrome sequence, using an approach analogous to the Euclidean algorithm for polynomials. The algorithm maintains a current estimate $\Lambda^{(r)}(x)$ at iteration $r$, and updates it based on the discrepancy between the predicted and actual next syndrome value.

\subsection{Concatenated Codes}
\label{sec:ec:concat}

\begin{theorem}[Forney 1966]
For any rate $R < C$, there exists a concatenated code with polynomial-time encoding/decoding and error probability $\exp(-n^\alpha)$.
\end{theorem}

Outer code: RS over $\mathbb{F}_{2^k}$ (corrects symbol errors). Inner code: binary code (e.g., convolutional) matched to channel. Inner decoder outputs symbols + reliability $\to$ outer decoder corrects remaining errors.

This construction was used in Voyager (1977) and deep space communications for decades.

Forney's key insight was \emph{divide and conquer}: rather than designing a single, monolithic code that simultaneously handles all noise patterns, split the problem into two layers. The inner code operates at the physical channel level, handling the raw bit errors. The outer code operates at a higher level, treating inner decoder failures as symbol errors. The overall error probability is the product of the two layers' error probabilities, leading to the exponential decay. Forney showed that with $n$ total blocklength, the error probability decays as $\exp(-n^\alpha)$ for some $\alpha > 0$, a dramatic improvement over uncoded transmission.

The Voyager spacecraft, launched in 1977, used a concatenated code: a $(255, 223)$ Reed-Solomon outer code over $\mathbb{F}_{256}$ concatenated with a $(2,1,7)$ convolutional inner code decoded by the Viterbi algorithm. This combination enabled communication from the edge of the solar system at data rates of just a few kilobits per second. The error correction was so effective that the Golden Record's images and sounds arrived at Earth with virtually no errors despite the 15-billion-kilometer journey.

\subsection{Convolutional Codes and Viterbi Algorithm}
\label{sec:ec:conv}

Convolutional code: infinite sequence, encoded by finite-state machine (shift register). Trellis diagram. Viterbi algorithm (1967): dynamic programming on trellis $\to$ ML sequence estimation in $O(n 2^\nu)$ where $\nu$ = constraint length.

\begin{example}[Convolutional Code and Viterbi Decoding: Worked Example]
Consider a rate $1/2$ convolutional code with constraint length $\nu = 2$ (memory $m = 2$). The encoder uses generators $g_1 = (1,1,1)$ and $g_2 = (1,0,1)$, producing two output bits per input bit. The shift register has 2 memory elements, giving $2^{m} = 4$ states: $00, 01, 10, 11$.

The trellis diagram shows state transitions:
\[
\begin{array}{c|c|c}
\text{Current state} & \text{Next state (input 0/1)} & \text{Output (input 0/1)} \\ \hline
00 & 00 / 10 & 00 / 11 \\
01 & 00 / 10 & 10 / 01 \\
10 & 01 / 11 & 11 / 00 \\
11 & 01 / 11 & 01 / 10
\end{array}
\]

Encode message $m = (1, 0, 1, 1, 0, 0, 0)$ (with 3 trailing zeros to flush the register):
\begin{itemize}
\item Start state: $00$.
\item Input $1 \to$ state $10$, output $11$.
\item Input $0 \to$ state $01$, output $10$.
\item Input $1 \to$ state $10$, output $00$.
\item Input $1 \to$ state $11$, output $01$.
\item Input $0 \to$ state $01$, output $10$.
\item Input $0 \to$ state $00$, output $11$.
\item Input $0 \to$ state $00$, output $00$.
\end{itemize}

Encoded sequence: $11\ 10\ 00\ 01\ 10\ 11\ 00$.

Now suppose the received sequence (with noise) is: $10\ 11\ 01\ 01\ 11\ 11\ 00$.

The Viterbi algorithm proceeds as follows:

\textbf{Step 1:} For each received pair, compute the Hamming distance to each possible output pair for every state transition.

\textbf{Step 2:} For each state at time $t$, maintain the minimum accumulated distance (path metric) and the surviving path.

\textbf{Step 3:} After processing all received symbols, trace back from the state with the minimum path metric.

For the first received pair $10$:
\begin{itemize}
\item State $00$ reached from $00$ (output $00$) with distance $d(10,00)=1$, or from $01$ (output $10$) with distance $d(10,10)=0$.
\item State $10$ reached from $00$ (output $11$) with distance $d(10,11)=1$, or from $01$ (output $01$) with distance $d(10,01)=1$.
\end{itemize}

Continuing through all received symbols, the algorithm maintains $4$ survivor paths. At the end, the path with minimum metric is traced back to recover the most likely transmitted sequence. The Viterbi algorithm is guaranteed to find the maximum-likelihood sequence — it is optimal but with complexity linear in $n$ and exponential only in constraint length $\nu$.

In practice, constraint lengths of $\nu = 7$ to $9$ are common, giving $2^7=128$ to $2^9=512$ states — easily manageable. For $\nu = 7$ on the AWGN channel, a rate $1/2$ convolutional code achieves about $5$ dB of coding gain at $10^{-5}$ BER.
\end{example}

\subsection{Turbo Codes}
\label{sec:ec:turbo}

Parallel concatenation of two convolutional codes with interleaver. Iterative decoding: two soft-in/soft-out (SISO) decoders exchange \emph{extrinsic information} (log-likelihood ratios). ``Turbo'' = feedback loop like turbocharger.

\begin{example}[Turbo Code Encoding and Decoding: Worked Example]
Consider a turbo encoder with two identical recursive systematic convolutional (RSC) component encoders, each of rate $1/2$ and constraint length $\nu = 3$.

\textbf{Encoder structure:}
\begin{itemize}
\item Input: information bits $u_1, u_2, \dots, u_k$.
\item First output (systematic): $u_1, u_2, \dots, u_k$ (the original bits).
\item Second output (parity 1): $p^{(1)}_1, p^{(1)}_2, \dots, p^{(1)}_k$ from the first RSC encoder.
\item Third output (parity 2): $p^{(2)}_1, p^{(2)}_2, \dots, p^{(2)}_k$ from the second RSC encoder, which receives the interleaved version of the input bits.
\end{itemize}

The overall rate is $1/3$, but puncturing can increase it to $1/2$ by alternately transmitting parity bits from each encoder.

\textbf{Iterative decoding:}
The turbo decoder consists of two SISO decoders connected in a loop. Each decoder computes log-likelihood ratios (LLRs):
\[
L(u_i) = \ln \frac{P(u_i = 1 \mid r)}{P(u_i = 0 \mid r)}
\]

The LLR decomposes into three components:
\[
L(u_i) = L_c \cdot r_i + L_a(u_i) + L_e(u_i)
\]
where $L_c \cdot r_i$ is the channel observation, $L_a(u_i)$ is the \emph{a priori} information (from the other decoder), and $L_e(u_i)$ is the \emph{extrinsic} information (new information generated by this decoder).

\textbf{Decoding iteration:}
\begin{enumerate}
\item Decoder 1 receives channel LLRs and parity $p^{(1)}$. Using $L_a(u_i) = 0$ initially, it computes $L_e^{(1)}(u_i) = L(u_i) - L_c \cdot r_i$ and passes this to decoder 2.
\item Decoder 2 interleaves $L_e^{(1)}$ to use as $L_a^{(2)}$, receives channel LLRs and parity $p^{(2)}$, computes $L_e^{(2)}$, de-interleaves it, and passes back to decoder 1.
\item Repeat for $6$--$12$ iterations. The extrinsic information converges: correct bits approach $+\infty$ (certain 1) and incorrect bits approach $-\infty$ (certain 0).
\end{enumerate}

After sufficient iterations, make hard decisions: $\hat{u}_i = 1$ if $L(u_i) > 0$, else $0$. Turbo codes achieve BER $< 10^{-5}$ at $E_b/N_0$ within 1 dB of the Shannon limit for large block sizes.

\textbf{Example with short block:} For a block of 4 bits with received LLRs $(0.5, -0.3, 1.2, -0.8)$:
\begin{itemize}
\item Iteration 1: Decoder 1 produces extrinisic values $(0.2, 0.1, -0.3, 0.4)$.
\item Decoder 2 uses these as priors, adds channel info, produces new extrinisic $(-0.1, 0.3, 0.5, -0.2)$.
\item Iteration 2: LLRs converge to $(0.6, 0.2, 1.4, -0.6)$ $\to$ decisions $(1, 1, 1, 0)$.
\end{itemize}

The critical insight: extrinsic information is the \emph{new} information about each bit that comes from the code structure, not from the channel or from the other decoder's previous output. By exchanging only extrinsic information, the loop avoids positive feedback (the turbo decoder would ``burn out'' if it recycled its own output).
\end{example}

\subsection{LDPC Codes}
\label{sec:ec:ldpc}

Sparse parity-check matrix $H$ (few 1s per row/col). Tanner graph: variable nodes (bits) + check nodes (parities). Belief propagation (sum-product algorithm) on graph. Iterative, local, parallelizable. Capacity-achieving for large blocklengths.

\begin{example}[LDPC Code: Small Example and Belief Propagation]
Consider a $(6, 3)$ LDPC code with parity-check matrix:
\[
H = \begin{pmatrix}
1 & 1 & 0 & 1 & 0 & 0 \\
0 & 1 & 1 & 0 & 1 & 0 \\
1 & 0 & 1 & 0 & 0 & 1
\end{pmatrix}
\]

The Tanner graph has 6 variable nodes $v_1, \dots, v_6$ (one per column) and 3 check nodes $c_1, c_2, c_3$ (one per row). An edge connects $v_j$ to $c_i$ if $H_{ij} = 1$.

Now consider transmission over the binary erasure channel (BEC) with erasure probability $\epsilon = 0.3$. The transmitted codeword is $c = (1, 0, 1, 1, 0, 1)$ and the received word is $r = (1, ?, 1, ?, 0, 1)$, where ``?'' denotes an erasure at positions 2 and 4.

\textbf{Belief propagation on the BEC:}
Each check node $c_i$ performs a parity check on its adjacent variable nodes. If all but one adjacent variable are known, the unknown value is determined by parity.

\textbf{Iteration 1:}
\begin{itemize}
\item $c_1$ connected to $v_1, v_2, v_4$. Known: $v_1 = 1, v_4 = ?$. Cannot resolve — two unknowns.
\item $c_2$ connected to $v_2, v_3, v_5$. Known: $v_3 = 1, v_5 = 0$. The check equation: $v_2 + v_3 + v_5 = 0$ (mod 2). So $v_2 + 1 + 0 = 0 \implies v_2 = 1$. Resolved!
\item $c_3$ connected to $v_1, v_3, v_6$. Known: $v_1 = 1, v_3 = 1, v_6 = 1$. Check equation: $1 + 1 + 1 = 1 \neq 0$. But $c$ is a codeword so $v_6 = 1$ is the correct value from transmission.
\end{itemize}

\textbf{Iteration 2:}
Now $v_2$ is known ($=1$).
\begin{itemize}
\item $c_1$ connected to $v_1, v_2, v_4$. Known: $v_1 = 1, v_2 = 1$. Check: $1 + 1 + v_4 = 0 \implies v_4 = 0$. Resolved!
\end{itemize}

After 2 iterations, all erasures are recovered. The decoding succeeds.

For the AWGN channel, belief propagation uses log-likelihood ratios instead of erasure symbols. The algorithm iterates:
\begin{enumerate}
\item Variable nodes send their LLRs to adjacent check nodes.
\item Check nodes compute updated LLRs using the $\tanh$ rule:
\[
L_{c \to v} = 2 \tanh^{-1}\left( \prod_{v' \in N(c) \setminus \{v\}} \tanh\left( \frac{L_{v' \to c}}{2} \right) \right)
\]
\item Variable nodes sum incoming LLRs with the channel LLR and send back.
\item After convergence, hard decisions from total LLRs.
\end{enumerate}

For large blocklengths ($n \approx 10^5$), LDPC codes with optimized degree distributions operate within 0.1 dB of the Shannon limit.
\end{example}

\subsubsection*{Density Evolution (Richardson-Urbanke 2001)}
Tracks the distribution of messages in BP decoding as blocklength $\to \infty$. For a given LDPC ensemble (degree distributions $\lambda(x), \rho(x)$), computes the threshold $p^*$ such that decoding succeeds iff channel error prob $< p^*$. Enables code design by optimization.

Density evolution analyzes the convergence of belief propagation in the asymptotic regime where the blocklength $N \to \infty$ and the Tanner graph is locally tree-like (no cycles up to an arbitrarily large neighborhood). For a $(d_l, d_r)$-regular LDPC ensemble, density evolution tracks the probability density of messages passed between variable and check nodes.

Let $P^{(l)}$ be the density of messages from variable nodes to check nodes at iteration $l$, and $Q^{(l)}$ the density of messages from check nodes to variable nodes. The update rules are:
\[
Q^{(l)} = \mathcal{C}^{-1}\left( \left[ \mathcal{C}(P^{(l-1)}) \right]^{\otimes (d_r - 1)} \right)
\]
\[
P^{(l)} = P_0 \otimes \left( Q^{(l)} \right)^{\otimes (d_l - 1)}
\]
where $\mathcal{C}$ denotes the cosh-Fourier transform, $\otimes$ is convolution, and $P_0$ is the channel message density. The decoding threshold $p^*$ is the supremum of channel parameters for which $P^{(l)} \to \delta_0$ (the zero-error density) as $l \to \infty$.

For the BEC, density evolution simplifies dramatically. The fraction of erasure messages at iteration $l$ follows:
\[
x_{l} = \epsilon \cdot \lambda(1 - \rho(1 - x_{l-1}))
\]
where $\lambda(x) = \sum_i \lambda_i x^{i-1}$ and $\rho(x) = \sum_i \rho_i x^{i-1}$ are the degree distributions from the edge perspective. The threshold $\epsilon^*$ is the largest $\epsilon$ for which $x_l \to 0$ as $l \to \infty$. For the $(3,6)$-regular ensemble, $\epsilon^* \approx 0.429$, compared to the Shannon limit of $0.5$ for the BEC.

Irregular LDPC codes, with carefully optimized degree distributions, achieve thresholds arbitrarily close to capacity. The design problem reduces to solving a linear programming optimization over degree distributions subject to the stability condition $\lambda'(0)\rho'(1) < 1/\epsilon$.

\subsubsection*{EXIT Charts (Ten Brink 2001)}
Extrinsic Information Transfer charts visualize the mutual information flow between component decoders in iterative decoding. Used to predict convergence and design code ensembles.

EXIT charts plot the \emph{extrinsic} mutual information $I_E$ at the output of one decoder against the \emph{a priori} mutual information $I_A$ at its input. For a given code and SNR, the EXIT function $T(I_A) = I_E$ describes how much new information the decoder produces.

For a turbo decoder with two component decoders, the EXIT chart shows the convergence trajectory:
\begin{itemize}
\item Decoder 1: $I_{E1} = T_1(I_{A1})$, where $I_{A1} = I_{E2}$ (interleaved).
\item Decoder 2: $I_{E2} = T_2(I_{A2})$, where $I_{A2} = I_{E1}$ (de-interleaved).
\end{itemize}

The decoding trajectory is a staircase between the two EXIT curves. If the curves intersect before reaching $(1,1)$, decoding gets stuck at a fixed point. If they remain separated up to $(1,1)$, decoding converges to vanishing error probability.

The key design rule: the area under the EXIT curve equals the code rate, and for capacity-achieving codes, the EXIT curves should be matched (decoder 1's curve is the mirror image of decoder 2's). EXIT charts thus provide a visual, intuitive tool for code design that complements the analytical rigor of density evolution.

\subsection{Polar Codes}
\label{sec:ec:polar}

Channel polarization: $W_N = W^{\otimes N} \to$ combine via $G_N = B_N F^{\otimes N}$ ($F = \begin{smallmatrix}1&0\\1&1\end{smallmatrix}$). Synthetic channels $W_N^{(i)}$ polarize: $I(W_N^{(i)}) \to 0$ or $1$. Choose indices with $I \approx 1$ for data; freeze others. Successive cancellation decoding $O(N \log N)$. List decoding improves finite-length performance.

\begin{example}[Channel Polarization for $N=2$: Detailed Example]
Start with the binary erasure channel $W$ with erasure probability $\epsilon = 0.5$. The channel $W$ has capacity $I(W) = 1 - \epsilon = 0.5$.

For $N = 2$, combine two independent copies of $W$ via the kernel $F = \begin{pmatrix}1 & 0\\ 1 & 1\end{pmatrix}$. The combined channel $W_2$ maps $(u_1, u_2)$ to $(y_1, y_2)$ through:
\[
(x_1, x_2) = (u_1 \oplus u_2, u_2)
\]
where $x_1, x_2$ are transmitted over independent copies of $W$, and $y_1, y_2$ are received.

Split $W_2$ into two synthetic channels:
\begin{itemize}
\item $W_2^{-}(y_1,y_2 \mid u_1)$: decode $u_1$ first (treat $u_2$ as noise).
\item $W_2^{+}(y_1,y_2,u_1 \mid u_2)$: decode $u_2$ after $u_1$ is known.
\end{itemize}

For the BEC, these synthetic channels are themselves BECs with erasure probabilities:
\[
\epsilon^{-} = 2\epsilon - \epsilon^2 = 2(0.5) - 0.25 = 0.75
\]
\[
\epsilon^{+} = \epsilon^2 = 0.25
\]

The capacities are:
\[
I(W_2^{-}) = 1 - 0.75 = 0.25,\quad
I(W_2^{+}) = 1 - 0.25 = 0.75
\]

Notice: $I(W_2^{-}) + I(W_2^{+}) = 1.0 = 2I(W)$, and $I(W_2^{-}) < I(W) < I(W_2^{+})$. The channels have polarized: one worse than the original, one better.

For $N = 4$, apply the same transformation recursively to each pair of synthetic channels. The erasure probabilities become:
\[
\begin{aligned}
\epsilon^{(1)} &= 2\epsilon^{-} - (\epsilon^{-})^2 = 2(0.75) - 0.5625 = 0.9375 \\
\epsilon^{(2)} &= (\epsilon^{-})^2 = 0.5625 \\
\epsilon^{(3)} &= 2\epsilon^{+} - (\epsilon^{+})^2 = 2(0.25) - 0.0625 = 0.4375 \\
\epsilon^{(4)} &= (\epsilon^{+})^2 = 0.0625
\end{aligned}
\]

For $N = 8$, the erasure probabilities are:
\[
(0.996, 0.879, 0.809, 0.316, 0.684, 0.191, 0.121, 0.004)
\]

The pattern is clear: as $N$ grows, the erasure probabilities approach either 0 or 1. The fraction of channels with $\epsilon \approx 0$ approaches $I(W) = 0.5$.

To achieve rate $R = 0.5$ on a BEC with $\epsilon = 0.5$, we:
\begin{itemize}
\item Compute $\epsilon^{(i)}$ for $i = 1, \dots, N$ (with $N$ large, say $1024$).
\item Sort by reliability. Select the $N/2$ channels with smallest $\epsilon^{(i)}$ for data.
\item Set frozen bits (known to both encoder and decoder, typically all-zeros) on the remaining $N/2$ channels.
\end{itemize}

For this example with $N = 8$ and rate $R = 0.5$, the information set is $\mathcal{A} = \{4, 5, 6, 7, 8\}$ (5 channels wait — actually we need exactly 4 for rate 1/2). The 4 best channels are indices 7, 8, 6, 4 with erasure probabilities $(0.121, 0.004, 0.191, 0.316)$. The remaining 4 channels are frozen.
\end{example}

\subsubsection*{Successive Cancellation Decoding}

Successive cancellation (SC) decoding estimates bits $u_1, u_2, \dots, u_N$ in order, using previous estimates as known values. For each bit $i$:

\[
\hat{u}_i = \begin{cases}
u_i \text{ (known frozen bit)} & \text{if } i \in \mathcal{A}^c \\
\arg\max_{u \in \{0,1\}} W_N^{(i)}(y_1^N, \hat{u}_1^{i-1} \mid u) & \text{if } i \in \mathcal{A}
\end{cases}
\]

The likelihood ratio for $u_i$ is computed recursively using the formulas:
\[
L_N^{(2i-1)}(y_1^N, \hat{u}_1^{2i-2}) = \frac{L_{N/2}^{(i)}(y_1^{N/2}, \hat{u}_{1,o}^{2i-2} \oplus \hat{u}_{1,e}^{2i-2}) \cdot L_{N/2}^{(i)}(y_{N/2+1}^N, \hat{u}_{1,e}^{2i-2}) + 1}{L_{N/2}^{(i)}(y_1^{N/2}, \hat{u}_{1,o}^{2i-2} \oplus \hat{u}_{1,e}^{2i-2}) + L_{N/2}^{(i)}(y_{N/2+1}^N, \hat{u}_{1,e}^{2i-2})}
\]
\[
L_N^{(2i)}(y_1^N, \hat{u}_1^{2i-1}) = \left[ L_{N/2}^{(i)}(y_1^{N/2}, \hat{u}_{1,o}^{2i-2} \oplus \hat{u}_{1,e}^{2i-2}) \right]^{1 - 2\hat{u}_{2i-1}} \cdot L_{N/2}^{(i)}(y_{N/2+1}^N, \hat{u}_{1,e}^{2i-2})
\]

where $\hat{u}_{1,o}$ and $\hat{u}_{1,e}$ are the odd and even indices of $\hat{u}_1^{2i-2}$, respectively.

The recursion halves the blocklength at each step, giving total complexity $O(N \log N)$. The decoding can be visualized as a binary tree traversal: root represents the full block, leaves are individual bits, and the algorithm performs a depth-first search computing LLRs bottom-up.

For the BEC, SC decoding is particularly simple. The decoder uses the same erasure probabilities computed by the Bhattacharyya parameter recursion:
\[
\beta^{(2i-1)} = 2\beta^{(i)} - (\beta^{(i)})^2,\quad
\beta^{(2i)} = (\beta^{(i)})^2
\]
where $\beta^{(i)}$ is the Bhattacharyya parameter of synthetic channel $W_N^{(i)}$.

\subsubsection*{Channel Polarization Proof Sketch}
For a B-DMC $W$, define two synthetic channels:
\begin{itemize}
\item $W^-(y_1,y_2|u_1) = \frac{1}{2} \sum_{u_2} W(y_1|u_1 \oplus u_2) W(y_2|u_2)$
\item $W^+(y_1,y_2,u_1|u_2) = \frac{1}{2} W(y_1|u_1 \oplus u_2) W(y_2|u_2)$
\end{itemize}
Then $I(W^-) + I(W^+) = 2I(W)$ and $I(W^-) \le I(W) \le I(W^+)$. Recursion $\to$ polarization.

The polarization phenomenon is proved by showing that the process of recursive channel combining and splitting is a martingale that converges almost surely to $0$ or $1$. The fraction of channels that become perfect equals the symmetric capacity $I(W)$. This is the Bhattacharyya process: a binary random tree where each leaf's Bhattacharyya parameter evolves according to the recursion above.

\textbf{Bhattacharyya parameter recursion:} For any B-DMC $W$ with Bhattacharyya parameter $Z(W) = \sum_y \sqrt{W(y|0)W(y|1)}$, the recursion gives:
\[
\beta^{-} = 2\beta - \beta^2,\quad \beta^{+} = \beta^2
\]
where $\beta = Z(W)$, $\beta^{-} = Z(W^-)$, and $\beta^{+} = Z(W^+)$. Since $Z(W) \approx 0$ for good channels and $Z(W) \approx 1$ for bad channels, the recursion provides an efficient way to evaluate channel quality without computing capacities explicitly.

\subsubsection*{List Decoding for Polar Codes}
Successive cancellation (SC) decoding makes hard decisions. SC-List (SCL) keeps $L$ best paths, uses CRC to select final codeword. Performance within $<0.5$ dB of capacity at finite lengths. Used in 5G control channel.

\subsubsection*{Polar Codes for Non-Binary Channels}
Polarization extends to $q$-ary channels and channels with memory (via kernel matrices larger than $F$). Recent work: $2 \times 2$ kernels for BSC, $16 \times 16$ kernels for faster polarization.

%======================================================================
% SECTION 6: CONSEQUENCES
%======================================================================
\section{Consequences: What Became Possible}
\label{sec:ec:consequences}

\begin{enumerate}
\item \textbf{Deep Space Communication:} Voyager (1977) used concatenated RS + convolutional. Later missions (Mars rovers, New Horizons) use Turbo/LDPC. Without coding, signals from billions of km would be unrecoverable.

\item \textbf{Consumer Storage:} CDs (CIRC: RS + interleaving), DVDs, Blu-ray, hard drives (LDPC), SSDs (LDPC/BCH), QR codes (RS). A CD with 1mm scratch plays perfectly.

\item \textbf{Wireless:} WiFi (LDPC since 802.11n), 4G/5G (Turbo/LDPC/Polar), Bluetooth (FEC). Spectral efficiency gains of 3--6$\times$ over uncoded.

\item \textbf{Internet:} TCP retransmits (ARQ) $\neq$ FEC. But QUIC, HTTP/3, satellite internet (Starlink) use FEC for latency reduction.

\item \textbf{RAID:} RAID-5/6 use parity/RS for disk failure tolerance. Erasure coding (RS, fountain codes) in distributed storage (Hadoop, Ceph, cloud).

\item \textbf{DNA Storage:} Synthesis/sequencing errors $\to$ need codes for substitution/indel channels. New frontier for coding theory.

\item \textbf{Quantum Error Correction:} Shor code (1995), Surface codes — coding theory for qubits. The threshold theorem: if physical error rate $< 1\%$, arbitrarily long quantum computation is possible.

\item \textbf{Economic Impact:} The coding gain from advanced error correction translates directly into cost savings. A 3 dB coding gain means you can halve the transmitter power, double the range, or double the data rate for the same power. In satellite communications, this is worth millions of dollars per dB. In cellular networks, turbo and LDPC codes enabled the transition from voice-only 2G to broadband 4G/5G by squeezing more bits per Hz from the limited spectrum.

\item \textbf{Space Exploration Beyond Voyager:} The Cassini mission to Saturn, the Mars Exploration Rovers, and the New Horizons flyby of Pluto all used increasingly sophisticated codes. The Mars rovers achieved data rates of 2--6 Mbps from Mars using LDPC codes — a dramatic improvement over the few kbps of earlier missions. Without coding, the stunning images from Titan, Mars, and Pluto would never have reached Earth.

\item \textbf{Optical Fiber Communication:} Modern fiber optic systems use powerful LDPC codes with 15--25\% overhead to achieve data rates of 400 Gbps and beyond per wavelength. The coding gain of 8--10 dB enables transoceanic cables that span 10,000+ km without regeneration.
\end{enumerate}

%======================================================================
% SECTION 7: PHILOSOPHICAL SIGNIFICANCE
%======================================================================
\section{Philosophical Significance: What It Revealed About Limits of Formal Systems}
\label{sec:ec:philosophical}

\begin{enumerate}
\item \textbf{Redundancy Is Not Waste — It's Structure:} The pre-Shannon view: redundancy = inefficiency. Coding theory: redundancy = \emph{controlled structure} that enables recovery. The right redundancy is \emph{information about the information}.

\item \textbf{Reliability from Unreliable Parts:} Von Neumann (1956): reliable computation from unreliable components (probabilistic logic). Error correction is the same principle: a noisy channel + a code = a noiseless channel. \emph{Composition of unreliable parts yields reliable wholes.}

\item \textbf{The Graph Perspective:} LDPC/Turbo revealed that codes are \emph{graphical models}. Decoding = inference on a graph. This unified coding theory with statistical physics (spin glasses), belief propagation, and machine learning (Boltzmann machines, variational inference).

\item \textbf{Deterministic vs. Random:} Shannon used random codes (existence). Practical codes are deterministic/structured. Polar codes bridged the gap: deterministic construction, provable capacity-achieving. The tension between \emph{randomness as proof technique} and \emph{structure as engineering requirement} drives the field.

\item \textbf{The Bit Is Physical (Landauer) $\to$ The Bit Is Redundant (Shannon):} Landauer (1961): erasing a bit costs $kT \ln 2$ energy. Shannon: transmitting a bit reliably costs \emph{extra bits}. The physical cost of reliability is redundancy.

\item \textbf{The Universality of Coding:} Error-correcting codes have appeared in unexpected places. The genetic code uses redundancy (triplet codons) to protect against mutations. Neural representations use population coding (distributed redundancy) for robustness. The principle — reliable function from unreliable components through structured redundancy — appears to be a universal design pattern in both natural and engineered systems.
\end{enumerate}

%======================================================================
% SECTION 8: MODERN LEGACY
%======================================================================
\section{Modern Legacy: Where This Idea Appears Today}
\label{sec:ec:legacy}

\begin{itemize}
\item \textbf{5G NR:} LDPC (data), Polar (control).
\item \textbf{WiFi 6/7:} LDPC.
\item \textbf{Starlink:} LDPC + interleaving.
\item \textbf{SSDs:} LDPC with soft-decision decoding (multiple reads at different voltages).
\item \textbf{Distributed Storage:} Fountain codes (LT, Raptor) — rateless, ideal for multicast.
\item \textbf{Quantum Error Correction:} Shor code (1995), Surface codes — coding theory for qubits.
\item \textbf{Machine Learning:} Neural decoders (learning BP schedules), code design via RL, transformers for decoding.
\item \textbf{DNA Storage:} Codes for insertion/deletion/substitution channels (Levenshtein, Helberg, etc.).
\item \textbf{6G and Beyond:} Polar codes for URLLC, LDPC for eMBB, joint source-channel coding, semantic communication.
\item \textbf{Optical Networks:} 400G/800G Ethernet uses soft-decision FEC with LDPC codes. The ITU-T G.709 standard specifies concatenated FEC frameworks for optical transport networks.
\item \textbf{Automotive:} CAN-FD, automotive Ethernet, and V2X communication all employ error correction for safety-critical applications. The ISO 26262 functional safety standard mandates coding for automotive electronics.
\item \textbf{Space-Solaris and Deep Space Network:} NASA's Deep Space Network now uses rate-1/2 and rate-2/3 LDPC codes standardized by CCSDS. These codes enable communication with spacecraft at Jupiter, Saturn, and beyond.
\end{itemize}

\begin{itemize}
\item Chapter 5 (Information): Channel capacity is the \emph{limit}; codes are the \emph{construction}.
\item Chapter 7 (Kolmogorov): Error correction $\leftrightarrow$ compressibility. A correctable error pattern has low Kolmogorov complexity.
\item Chapter 10 (Cryptography): Error correction $\neq$ encryption, but both use finite fields. RS codes in secret sharing (Shamir).
\item Chapter 13 (Consensus): Error correction in Byzantine agreement (coded redundancy for fault tolerance).
\end{itemize}

%======================================================================
% SECTION 9: LESSONS FOR FUTURE SCIENTISTS
%======================================================================
\section{Summary}
\label{sec:ec:summary}

This chapter surveyed the theory and practice of error-correcting codes, tracing the development from Hamming codes (1950) to modern capacity-approaching constructions. The chapter began with algebraic codes: Hamming codes were presented as the first systematic construction of single-error-correcting codes, with the $[7,4,3]$ code achieving perfect sphere packing. Reed--Solomon codes over $\mathbb{F}_q$ were introduced as maximum distance separable (MDS) codes with minimum distance $d = n-k+1$, optimal for correcting burst errors. BCH codes and their algebraic decoding via the Berlekamp--Massey algorithm were discussed. Forney's concatenated codes (1966) demonstrated that algebraic codes could approach the Shannon limit when combined with interleaving. Convolutional codes and the Viterbi algorithm (1967) provided a trellis-based maximum-likelihood decoding framework. The chapter then examined turbo codes (Berrou, Glavieux, Thitimajshima 1993), which achieved within 0.5 dB of capacity using parallel concatenated convolutional codes and iterative belief propagation decoding. Low-density parity-check (LDPC) codes (Gallager 1963) were revisited through the lens of density evolution (Richardson, Urbanke, Shokrollahi 1999--2001), which characterized their decoding thresholds on the binary erasure channel and binary symmetric channel. Polar codes (Arıkan 2009) were presented as the first provably capacity-achieving code with polynomial-time encoding and decoding, based on the phenomenon of channel polarization. The chapter concluded with a discussion of finite-length performance, the gap between Shannon's random-coding bound and explicit constructions, and the theoretical lower bounds on blocklength for a given error probability (Polyanskiy, Poor, Verdú 2010).

\section*{Key Takeaways}
\begin{itemize}
\item Algebraic codes (Hamming, Reed--Solomon, BCH) provide explicit constructions with provable distance guarantees and efficient algebraic decoding algorithms.
\item Turbo codes and LDPC codes achieve near-capacity performance through iterative message-passing decoding on graphical models.
\item Polar codes are the first explicit construction provably achieving capacity for binary-input symmetric memoryless channels, via channel polarization.
\item The gap between Shannon's existential random-coding proof and explicit capacity-achieving constructions with efficient decoding spanned approximately 60 years.
\item Open problems include the finite-blocklength scaling law for polar and LDPC codes, the design of universal codes for unknown channel statistics, and quantum error-correcting codes approaching the quantum capacity.
\end{itemize}

The theory-to-practice pipeline for codes is brutally slow. After Gallager's 1963 thesis, LDPC codes waited 33 years for rediscovery and another 13 years for adoption in WiFi. Berrou's turbo codes went from ICC '93 to 3G in about 7 years --- considered lightning fast. The lesson: revolutionary ideas often face a generation of resistance before they become infrastructure.

%======================================================================
% SECTION 10: EXERCISES
%======================================================================
\section{Exercises}
\label{sec:ec:exercises}

\begin{exercise}[Basic]
Construct the generator and parity-check matrices for the $[7,4,3]$ Hamming code. Encode message $1011$. Introduce error in bit 3. Compute syndrome and correct.
\end{exercise}

\begin{exercise}[Basic]
Show that a Reed-Solomon $[n,k]$ code over $\mathbb{F}_q$ has minimum distance $n-k+1$ (MDS). Why does this make RS codes ideal for burst errors?
\end{exercise}

\begin{exercise}[Basic]
For the $(7,4)$ Hamming code, list all 16 codewords. Compute the minimum distance. Verify that the code is perfect by checking that spheres of radius 1 around each codeword partition the space $\mathbb{F}_2^7$.
\end{exercise}

\begin{exercise}[Basic]
Construct a syndrome decoding table for the $(7,4)$ Hamming code. Use it to decode the received word $r = (1,1,0,1,0,0,1)$. Show your steps. Then introduce a double-bit error and explain why the decoder makes an error.
\end{exercise}

\begin{exercise}[Intermediate]
Implement the Berlekamp-Massey algorithm for decoding BCH/RS codes. Decode a $(15,7)$ RS code over $\mathbb{F}_{16}$ with 4 symbol errors.
\end{exercise}

\begin{exercise}[Intermediate]
For a $(2,1,3)$ convolutional code with generators $g_1 = (1,1,1)$ and $g_2 = (1,0,1)$, draw the trellis diagram. Encode the message $m = (1,0,1,1)$. Decode the received sequence $r = (11, 10, 00, 01, 10)$ using the Viterbi algorithm, showing the survivor paths at each step.
\end{exercise}

\begin{exercise}[Intermediate]
Consider a turbo decoder with two identical RSC component codes. For a block of 4 information bits, the received LLRs are $(-0.2, 0.8, 1.1, -0.5)$. The first decoder produces extrinsic information $(0.3, -0.1, 0.4, 0.2)$. Compute the input LLRs for the second decoder and perform one decoding iteration. Show how the extrinsic information evolves.
\end{exercise}

\begin{exercise}[Intermediate]
For the $(6,3)$ LDPC code with parity-check matrix $H = [1\ 1\ 0\ 1\ 0\ 0;\ 0\ 1\ 1\ 0\ 1\ 0;\ 1\ 0\ 1\ 0\ 0\ 1]$, draw the Tanner graph. If the received word on the BEC is $r = (1, ?, 0, ?, 1, 0)$ where ``?'' denotes an erasure, run belief propagation to recover the erased bits.
\end{exercise}

\begin{exercise}[Intermediate]
For the binary erasure channel with $\epsilon = 0.4$, compute the Bhattacharyya parameters for $N = 8$ levels of polarization. Identify which channel indices should be used for data at rate $R = 0.5$. Which frozen bit pattern maximizes the code's performance?
\end{exercise}

\begin{exercise}[Challenge]
Research the \emph{channel polarization} proof for polar codes. Explain why $I(W^-)+I(W^+) = 2I(W)$ and $I(W^-) \le I(W) \le I(W^+)$. Implement the Bhattacharyya parameter recursion for $W = \text{BSC}(p)$.
\end{exercise}

\begin{exercise}[Challenge]
Research \emph{EXIT charts} for LDPC codes. Plot the EXIT functions for a $(3,6)$-regular LDPC ensemble and show how the area property relates to capacity.
\end{exercise}

\begin{exercise}[Challenge]
Derive the density evolution equations for a $(3,6)$-regular LDPC code on the BEC. Compute the decoding threshold $\epsilon^*$ and compare with the Shannon limit $1 - R = 0.5$. Explain the gap.
\end{exercise}

\begin{exercise}[Challenge]
Prove that the $(7,4)$ Hamming code is perfect using the sphere-packing bound. Generalize: for which parameters $r$ does a perfect Hamming code exist? Why are there no nontrivial perfect codes beyond Hamming and Golay?
\end{exercise}

\begin{exercise}[Challenge]
Implement Polar SC-List decoding with CRC. Compare BLER vs SNR for a (1024, 512) polar code with list sizes $L=1, 2, 4, 8, 32$.
\end{exercise}

\begin{exercise}[Computational]
Simulate a $(3,6)$-regular LDPC code on BSC($p=0.1$) using belief propagation. Plot BER vs. iterations. Compare with ML decoding (exhaustive for small $n$).
\end{exercise}

\begin{exercise}[Computational]
Implement a Polar code encoder/decoder. Use it to transmit an image over a simulated BSC. Compare PSNR with uncoded transmission and with a (7,4) Hamming code.
\end{exercise}

\begin{exercise}[Computational]
Simulate the Viterbi algorithm for the $(2,1,3)$ convolutional code on the AWGN channel. Plot BER vs. $E_b/N_0$ for constraint lengths $\nu = 3, 5, 7$. Compare the coding gain at BER $10^{-4}$.
\end{exercise}

\begin{exercise}[Computational]
Implement a Reed-Solomon encoder and decoder over $\mathbb{F}_{16}$ using a $(15, 11)$ RS code. Inject random symbol errors and measure the decoder's failure rate vs. the number of errors. Compare with the theoretical bound $t = \lfloor (n-k)/2 \rfloor$.
\end{exercise}

\begin{exercise}[Computational]
Compare the finite-length performance of $(7,4)$ Hamming, $(15,11)$ Hamming, $(1024,512)$ irregular LDPC, and $(1024,512)$ polar codes on the binary AWGN channel. Plot BER vs. $E_b/N_0$ for all four codes on the same axes. What blocklength is needed for each code to reach within 1 dB of capacity?
\end{exercise}

\begin{exercise}[Computational]
Implement the Berlekamp-Massey algorithm for the $(15,7)$ RS code over $\mathbb{F}_{16}$. Given a received word with 3 symbol errors, show the syndrome computation, the BM iterations, the Chien search, and Forney's formula step by step. Verify that the decoding succeeds.
\end{exercise}

%======================================================================
% TIMELINE
%======================================================================
\section{Timeline}
\label{sec:ec:timeline}
\addcontentsline{toc}{section}{Timeline}

\begin{tabularx}{\linewidth}{|l|X|}
\hline
\textbf{Year} & \textbf{Event} \\ \hline
1948 & Shannon: Channel Coding Theorem (random codes) \\ \hline
1950 & Hamming: (7,4) Hamming code; algebraic coding theory \\ \hline
1954 & Golay: Perfect binary (23,12,7) and ternary (11,6,5) codes \\ \hline
1959 & Bose, Chaudhuri, Hocquenghem: BCH codes \\ \hline
1960 & Reed, Solomon: Polynomial codes over $\mathbb{F}_q$ (RS codes) \\
 \hline
1960 & Peterson: BCH decoding algorithm \\
 \hline
1963 & Gallager: LDPC codes, sum-product algorithm (PhD thesis) \\
 \hline
1966 & Forney: Concatenated codes \\
 \hline
1967 & Viterbi: Convolutional code decoding (dynamic programming) \\
 \hline
1968 & Berlekamp: Algebraic decoding of BCH/RS (Massey 1969: BM algorithm) \\
 \hline
1977 & Voyager: Concatenated RS + convolutional code in deep space \\
 \hline
1982 & CD standard: CIRC (Cross-Interleaved RS Code) \\
 \hline
1993 & Berrou, Glavieux, Thitimajshima: Turbo codes (ICC '93) \\
 \hline
1996 & MacKay, Neal; Richardson, Urbanke: LDPC rediscovery, density evolution \\
 \hline
2001 & Shokrollahi: Raptor codes (fountain codes) \\
 \hline
2009 & Arıkan: Polar codes, channel polarization \\
 \hline
2016 & 5G NR adopts LDPC (data) + Polar (control) \\
 \hline
2022 & 3GPP Release 17: Enhanced polar codes for URLLC \\
 \hline
\end{tabularx}

%======================================================================
% CITATIONS
%======================================================================

%======================================================================
% INDEX ENTRIES
%======================================================================
\index{error-correcting codes}
\index{Hamming code}
\index{Reed-Solomon code}
\index{BCH code}
\index{concatenated codes}
\index{convolutional code}
\index{Viterbi algorithm}
\index{Turbo code}
\index{LDPC code}
\index{Polar code}
\index{channel polarization}
\index{belief propagation}
\index{iterative decoding}
\index{Gallager, Robert}
\index{Berrou, Claude}
\index{Arıkan, Erdal}
\index{Shannon, Claude}
\index{Hamming, Richard}
\index{density evolution}
\index{EXIT charts}
\index{list decoding}
\index{quantum error correction}
\index{surface codes}
\index{Shor code}
\index{Berlekamp-Massey algorithm}
\index{Chien search}
\index{Forney's formula}
\index{successive cancellation decoding}
\index{Bhattacharyya parameter}
\index{Tanner graph}
\index{extrinsic information}
\index{standard array}
\index{syndrome decoding}
\index{Forney, David}
\index{Voyager spacecraft}
\index{5G NR}
\index{Starlink}
\index{Reed, Irving}
\index{Solomon, Gustave}
